\chapter{The Ultimate Guide to Japanese Convenience Store}

% ================= INTRO =================
\begin{convobox}{1}
    \noindent
    \jpkana{みなさん}{Minasan} \jpkana{こんにちは、}{konnichiwa,} \jpkana{たなか}{Tanaka} \jpkana{です。}{desu.}

    \vspace{1.5em}
    \noindent{\Large \color{articolor} \textbf{Arti:} Halo semuanya, saya Tanaka.}
\end{convobox}

\begin{convobox}{2}
    \noindent
    \jpword{きょう}{今日}{Kyou} \jpkana{は、}{wa,} 
    \jpkana{コンビニ}{konbini} \jpkana{での}{de no} 
    \jpword{にほんご}{日本語}{Nihongo} \jpword{かいわ}{会話}{kaiwa} \jpkana{を}{o} 
    \jpword{いっしょ}{一緒}{issho} \jpkana{に}{ni} 
    \jpword{れんしゅう}{練習}{renshuu} \jpkana{しましょう。}{shimashou.}

    \vspace{1.5em}
    \noindent{\Large \color{articolor} \textbf{Arti:} Hari ini, mari kita berlatih percakapan bahasa Jepang di minimarket bersama-sama.}
\end{convobox}

\begin{convobox}{3}
    \noindent
    \jpkana{みなさん}{Minasan} \jpkana{が}{ga} 
    \jpword{にほん}{日本}{Nihon} \jpkana{に}{ni} \jpword{き}{来}{ki}\jpkana{たら、}{tara,} 
    \jpkana{きっと}{kitto} \jpword{いっかい}{1回}{ikkai} \jpkana{は}{wa} 
    \jpkana{コンビニ}{konbini} \jpkana{に}{ni} 
    \jpword{い}{行}{i}\jpkana{くと}{ku to} 
    \jpword{おも}{思}{omo}\jpkana{います。}{imasu.}

    \vspace{1.5em}
    \noindent{\Large \color{articolor} \textbf{Arti:} Jika kalian datang ke Jepang, aku yakin pasti setidaknya 1 kali akan pergi ke minimarket.}
\end{convobox}

\begin{convobox}{4}
    \noindent
    \jpkana{コンビニ}{Konbini} \jpkana{での}{de no} 
    \jpword{かいわ}{会話}{kaiwa} \jpkana{は、}{wa,} 
    \jpword{いぜん}{以前}{izen} \jpkana{の}{no} 
    \jpword{どうが}{動画}{douga} \jpkana{でも}{demo} 
    \jpword{しょうかい}{紹介}{shoukai} \jpkana{したことがありますが、}{shita koto ga arimasu ga,}

    \vspace{1.5em}
    \noindent{\Large \color{articolor} \textbf{Arti:} Percakapan di minimarket sudah pernah aku perkenalkan di video sebelumnya, tapi}
\end{convobox}

\begin{convobox}{5}
    \noindent
    \jpword{こんかい}{今回}{Konkai} \jpkana{の}{no} 
    \jpword{どうが}{動画}{douga} \jpkana{では、}{de wa,} 
    \jpkana{レジで}{reji de} 
    \jpword{てんいん}{店員}{tenin}\jpkana{さんに}{san ni} 
    \jpkana{よく}{yoku} 
    \jpword{き}{聞}{ki}\jpkana{かれる}{kareru} 
    \jpword{しつもん}{質問}{shitsumon} \jpkana{を}{o} 
    \jpword{ひと}{一}{hito}\jpkana{つずつ}{tsu zutsu} 
    \jpword{いっしょ}{一緒}{issho} \jpkana{に}{ni} 
    \jpword{かくにん}{確認}{kakunin} \jpkana{します。}{shimasu.}

    \vspace{1.5em}
    \noindent{\Large \color{articolor} \textbf{Arti:} Di video kali ini, kita akan mengecek satu per satu pertanyaan yang sering ditanyakan oleh kasir di mesin kasir.}
\end{convobox}

\begin{convobox}{6}
    \noindent
    \jpword{どうが}{動画}{Douga} \jpkana{の}{no} 
    \jpword{こうはん}{後半}{kouhan} \jpkana{では、}{de wa,} 
    \jpkana{ロールプレイ}{rooru purei} \jpkana{で}{de} 
    \jpword{じっせんてき}{実践的}{jissenteki} \jpkana{な}{na} 
    \jpword{かいわ}{会話}{kaiwa} \jpkana{の}{no} 
    \jpword{れんしゅう}{練習}{renshuu} \jpkana{も}{mo} \jpkana{するので}{suru node}

    \vspace{1.5em}
    \noindent{\Large \color{articolor} \textbf{Arti:} Di paruh kedua video, karena kita juga akan berlatih percakapan praktis dengan bermain peran (roleplay),}
\end{convobox}

\begin{convobox}{7}
    \noindent
    \jpkana{ぜひ}{Zehi} 
    \jpword{さいご}{最後}{saigo} \jpkana{まで}{made} 
    \jpword{み}{観}{mi}\jpkana{てください。}{te kudasai.}

    \vspace{1.5em}
    \noindent{\Large \color{articolor} \textbf{Arti:} pastikan untuk menonton sampai akhir ya.}
\end{convobox}

\begin{convobox}{8}
    \noindent
    \jpkana{それでは}{Sore dewa} 
    \jpword{さっそく}{早速}{sassoku} 
    \jpword{い}{行}{i}\jpkana{ってみましょう！}{tte mimashou!}

    \vspace{1.5em}
    \noindent{\Large \color{articolor} \textbf{Arti:} Kalau begitu, ayo langsung kita mulai!}
\end{convobox}


% ================= TITLE: CONVERSATION =================
\vspace{2em}
\begin{tcolorbox}[colback=boxframe, colframe=boxframe, rounded corners, boxrule=0pt, arc=3mm, left=2mm, top=2mm, bottom=2mm]
    \Large \bfseries \color{white} 🏪 1. Conversation (Percakapan Kecepatan Normal)
\end{tcolorbox}
\vspace{1em}

\begin{convobox}{9}
    \noindent
    \jpword{まず}{先ず}{Mazu} \jpword{いちど}{一度}{ichido} 
    \jpkana{コンビニ}{konbini} \jpkana{での}{de no} 
    \jpword{かいわ}{会話}{kaiwa} \jpkana{を}{o} 
    \jpword{しぜん}{自然}{shizen} \jpkana{な}{na} 
    \jpkana{スピードで}{supiido de} 
    \jpword{き}{聞}{ki}\jpkana{いてみましょう。}{ite mimashou.}

    \vspace{1.5em}
    \noindent{\Large \color{articolor} \textbf{Arti:} Pertama-tama, mari kita dengarkan sekali percakapan di minimarket dengan kecepatan natural.}
\end{convobox}

\begin{convobox}{10}
    \noindent
    \jpkana{いらっしゃいませ。}{Irasshaimase.} 
    \jpkana{こちら}{Kochira} 
    \jpword{あたた}{温}{atata}\jpkana{めますか？}{memasu ka?}

    \vspace{1.5em}
    \noindent{\Large \color{articolor} \textbf{Arti:} Selamat datang. Apakah ini mau dipanaskan?}
\end{convobox}

\begin{convobox}{11}
    \noindent
    \jpkana{いえ、}{Ie,} \jpword{だいじょうぶ}{大丈夫}{daijoubu} \jpkana{です。}{desu.}

    \vspace{1.5em}
    \noindent{\Large \color{articolor} \textbf{Arti:} Tidak, tidak usah/tidak apa-apa.}
\end{convobox}

\begin{convobox}{12}
    \noindent
    \jpkana{お}{O}\jpword{はし}{箸}{hashi} \jpkana{は}{wa} 
    \jpword{りよう}{利用}{riyou} \jpkana{ですか？}{desu ka?}

    \vspace{1.5em}
    \noindent{\Large \color{articolor} \textbf{Arti:} Apakah menggunakan (membutuhkan) sumpit?}
\end{convobox}

\begin{convobox}{13}
    \noindent
    \jpkana{はい、}{Hai,} 
    \jpkana{お}{o}\jpword{ねが}{願}{nega}\jpkana{いします。}{ishimasu.}

    \vspace{1.5em}
    \noindent{\Large \color{articolor} \textbf{Arti:} Ya, tolong (berikan).}
\end{convobox}

\begin{convobox}{14}
    \noindent
    \jpword{ふくろ}{袋}{Fukuro} \jpkana{は}{wa} 
    \jpword{りよう}{利用}{riyou} \jpkana{ですか？}{desu ka?}

    \vspace{1.5em}
    \noindent{\Large \color{articolor} \textbf{Arti:} Apakah menggunakan (membutuhkan) kantong plastik?}
\end{convobox}

\begin{convobox}{15}
    \noindent
    \jpkana{はい、}{Hai,} 
    \jpkana{お}{o}\jpword{ねが}{願}{nega}\jpkana{いします。}{ishimasu.}

    \vspace{1.5em}
    \noindent{\Large \color{articolor} \textbf{Arti:} Ya, tolong.}
\end{convobox}

\begin{convobox}{16}
    \noindent
    \jpword{いちまい}{1枚}{Ichi-mai} 
    \jpword{さんえん}{3円}{san-en} \jpkana{ですが、}{desu ga,} 
    \jpkana{よろしいでしょうか？}{yoroshii deshou ka?}

    \vspace{1.5em}
    \noindent{\Large \color{articolor} \textbf{Arti:} 1 lembarnya 3 yen, apakah tidak apa-apa?}
\end{convobox}

\begin{convobox}{17}
    \noindent
    \jpkana{はい。}{Hai.}

    \vspace{1.5em}
    \noindent{\Large \color{articolor} \textbf{Arti:} Ya.}
\end{convobox}

\begin{convobox}{18}
    \noindent
    \jpkana{ポイントカード}{Pointo kaado} \jpkana{は}{wa} 
    \jpkana{お}{o}\jpword{も}{持}{mo}\jpkana{ちですか？}{chi desu ka?}

    \vspace{1.5em}
    \noindent{\Large \color{articolor} \textbf{Arti:} Apakah Anda memiliki/membawa kartu poin?}
\end{convobox}

\begin{convobox}{19}
    \noindent
    \jpkana{いえ。}{Ie.}

    \vspace{1.5em}
    \noindent{\Large \color{articolor} \textbf{Arti:} Tidak.}
\end{convobox}

\begin{convobox}{20}
    \noindent
    \jpkana{お}{o}\jpword{かいけい}{会計}{kaikei} 
    \jpword{ろっぴゃくえん}{600円}{roppyaku-en} 
    \jpkana{でございます。}{de gozaimasu.}

    \vspace{1.5em}
    \noindent{\Large \color{articolor} \textbf{Arti:} Total pembayarannya 600 yen.}
\end{convobox}

\begin{convobox}{21}
    \noindent
    \jpkana{Tanaka Pay}{Tanaka Pay} \jpkana{で}{de} 
    \jpkana{お}{o}\jpword{ねが}{願}{nega}\jpkana{いします。}{ishimasu.}

    \vspace{1.5em}
    \noindent{\Large \color{articolor} \textbf{Arti:} Pakai Tanaka Pay, tolong.}
\end{convobox}

\begin{convobox}{22}
    \noindent
    \jpkana{かしこまりました。}{Kashikomarimashita.}

    \vspace{1.5em}
    \noindent{\Large \color{articolor} \textbf{Arti:} Baik, saya mengerti (Sangat Sopan).}
\end{convobox}

\begin{convobox}{23}
    \noindent
    \jpkana{レシート}{Reshiito} 
    \jpword{りよう}{利用}{riyou} \jpkana{ですか？}{desu ka?}

    \vspace{1.5em}
    \noindent{\Large \color{articolor} \textbf{Arti:} Apakah butuh struk?}
\end{convobox}

\begin{convobox}{24}
    \noindent
    \jpkana{いえ、}{Ie,} \jpword{だいじょうぶ}{大丈夫}{daijoubu} \jpkana{です。}{desu.}

    \vspace{1.5em}
    \noindent{\Large \color{articolor} \textbf{Arti:} Tidak, tidak usah.}
\end{convobox}

\begin{convobox}{25}
    \noindent
    \jpkana{ありがとうございました。}{Arigatou gozaimashita.} 
    \jpkana{また}{Mata} 
    \jpkana{お}{o}\jpword{こ}{越}{ko}\jpkana{しくださいませ。}{shi kudasaimase.}

    \vspace{1.5em}
    \noindent{\Large \color{articolor} \textbf{Arti:} Terima kasih banyak. Silakan datang kembali.}
\end{convobox}

\begin{convobox}{26}
    \noindent
    \jpkana{それでは、}{Sore dewa,} 
    \jpword{じゅんばん}{順番}{junban} \jpkana{に}{ni} 
    \jpword{せつめい}{説明}{setsumei} \jpkana{しますね。}{shimasu ne.}

    \vspace{1.5em}
    \noindent{\Large \color{articolor} \textbf{Arti:} Kalau begitu, akan saya jelaskan secara berurutan ya.}
\end{convobox}

% ================= TITLE: PHRASE 1 =================
\vspace{2em}
\begin{tcolorbox}[colback=boxframe, colframe=boxframe, rounded corners, boxrule=0pt, arc=3mm, left=2mm, top=2mm, bottom=2mm]
    \Large \bfseries \color{white} 🍱 2. Phrase 1 (Memanaskan Makanan)
\end{tcolorbox}
\vspace{1em}

\begin{convobox}{27}
    \noindent
    \jpkana{「こちら}{"Kochira} \jpword{あたた}{温}{atata}\jpkana{めますか？」}{memasu ka?"}

    \vspace{1.5em}
    \noindent{\Large \color{articolor} \textbf{Arti:} "Apakah ini mau dipanaskan?"}
\end{convobox}

\begin{convobox}{28}
    \noindent
    \jpword{べんとう}{弁当}{Bentou} \jpkana{や}{ya} 
    \jpkana{パスタ、}{pasuta,} 
    \jpkana{ラーメンなど、}{raamen nado,} 
    \jpword{あたた}{温}{atata}\jpkana{めることができる}{meru koto ga dekiru} 
    \jpword{しょうひん}{商品}{shouhin} \jpkana{は、}{wa,}

    \vspace{1.5em}
    \noindent{\Large \color{articolor} \textbf{Arti:} Produk yang bisa dipanaskan seperti kotak bekal (Bento), pasta, atau ramen,}
\end{convobox}

\begin{convobox}{29}
    \noindent
    \jpword{てんいん}{店員}{Tenin}\jpkana{さんが}{san ga} 
    \jpword{でんし}{電子}{denshi}\jpkana{レンジで}{renji de} 
    \jpword{あたた}{温}{atata}\jpkana{めてくれます。}{mete kuremasu.}

    \vspace{1.5em}
    \noindent{\Large \color{articolor} \textbf{Arti:} akan dipanaskan oleh kasir di dalam microwave.}
\end{convobox}

\begin{convobox}{30}
    \noindent
    \jpword{こた}{答}{Kota}\jpkana{えるときは、}{eru toki wa,} 
    \jpkana{このように}{kono you ni} 
    \jpword{い}{言}{i}\jpkana{います。}{imasu.}

    \vspace{1.5em}
    \noindent{\Large \color{articolor} \textbf{Arti:} Saat menjawab, ucapkanlah seperti ini.}
\end{convobox}

\begin{convobox}{31}
    \noindent
    \jpkana{「はい、}{"Hai,} \jpkana{お}{o}\jpword{ねが}{願}{nega}\jpkana{いします。」}{ishimasu."}

    \vspace{1.5em}
    \noindent{\Large \color{articolor} \textbf{Arti:} "Ya, tolong."}
\end{convobox}

\begin{convobox}{32}
    \noindent
    \jpkana{「いえ、}{"Ie,} \jpword{だいじょうぶ}{大丈夫}{daijoubu} \jpkana{です。」}{desu."}

    \vspace{1.5em}
    \noindent{\Large \color{articolor} \textbf{Arti:} "Tidak, tidak usah."}
\end{convobox}

\begin{convobox}{33}
    \noindent
    \jpword{てんいん}{店員}{Tenin}\jpkana{さんに}{san ni} 
    \jpword{なに}{何}{nani}\jpkana{か}{ka} 
    \jpword{き}{聞}{ki}\jpkana{かれた}{kareta} 
    \jpword{とき}{時}{toki}\jpkana{は、}{wa,}

    \vspace{1.5em}
    \noindent{\Large \color{articolor} \textbf{Arti:} Saat ditanya sesuatu oleh kasir,}
\end{convobox}

\begin{convobox}{34}
    \noindent
    \jpword{きほんてき}{基本的}{Kihonteki} \jpkana{に}{ni} 
    \jpkana{この}{kono} \jpword{ふた}{二}{futa}\jpkana{つを}{tsu o} 
    \jpword{つか}{使}{tsuka}\jpkana{って}{tte} 
    \jpword{こた}{答}{kota}\jpkana{えれば}{ereba} 
    \jpword{だいじょうぶ}{大丈夫}{daijoubu} \jpkana{です。}{desu.}

    \vspace{1.5em}
    \noindent{\Large \color{articolor} \textbf{Arti:} pada dasarnya kalau kamu menjawab menggunakan 2 frasa ini saja sudah cukup (aman).}
\end{convobox}

\begin{convobox}{35}
    \noindent
    \jpword{しつもん}{質問}{Shitsumon} \jpkana{を}{o} 
    \jpword{き}{聞}{ki}\jpkana{き}{ki}\jpword{と}{取}{to}\jpkana{れなかった}{renakatta} 
    \jpword{とき}{時}{toki}\jpkana{は、}{wa,} 
    \jpkana{このように}{kono you ni} 
    \jpword{い}{言}{i}\jpkana{いましょう。}{imashou.}

    \vspace{1.5em}
    \noindent{\Large \color{articolor} \textbf{Arti:} Saat kamu tidak bisa mendengar pertanyaannya dengan jelas, ucapkanlah seperti ini.}
\end{convobox}

\begin{convobox}{36}
    \noindent
    \jpkana{「すみません、}{"Sumimasen,} 
    \jpword{もういちど}{もう一度}{mou ichido} 
    \jpkana{お}{o}\jpword{ねが}{願}{nega}\jpkana{いします。」}{ishimasu."}

    \vspace{1.5em}
    \noindent{\Large \color{articolor} \textbf{Arti:} "Maaf, tolong ulangi sekali lagi."}
\end{convobox}

\begin{lessonbox}{Aksi Pahlawan Konbini: Jurus 2 Jawaban Sakti!}
    \textbf{Penjelasan:} Di Jepang, kasir bicara sangaaat cepat dan sopan (Keigo). Jangan panik! Kamu TIDAK HARUS mengerti setiap kata.
    \\
    \textbf{🎬 ACTION / PRAKTEK:} 
    \begin{enumerate}
        \item Jika kasir menunjuk kotak bento/makananmu, mereka pasti bertanya soal \textit{Microwave} (Atatamemasu ka?).
        \item Cukup anggukkan kepala sedikit dan ucapkan keras-keras: \textbf{"Hai, onegaishimasu!"} (Ya tolong). 
        \item Jika kamu membeli minuman dingin atau tidak mau dipanaskan, gelengkan kepala dan ucapkan: \textbf{"Ie, daijoubu desu!"} (Tidak usah).
    \end{enumerate}
    Dua jurus ini (\textit{Hai, onegaishimasu} dan \textit{Ie, daijoubu desu}) adalah tameng utamamu di Konbini!
\end{lessonbox}

% ================= TITLE: PHRASE 2 =================
\vspace{2em}
\begin{tcolorbox}[colback=boxframe, colframe=boxframe, rounded corners, boxrule=0pt, arc=3mm, left=2mm, top=2mm, bottom=2mm]
    \Large \bfseries \color{white} 🥢 3. Phrase 2 (Sumpit \& Alat Makan)
\end{tcolorbox}
\vspace{1em}

\begin{convobox}{37}
    \noindent
    \jpkana{「お}{"O}\jpword{はし}{箸}{hashi} \jpkana{は}{wa} 
    \jpword{りよう}{利用}{riyou} \jpkana{ですか？」}{desu ka?"}

    \vspace{1.5em}
    \noindent{\Large \color{articolor} \textbf{Arti:} "Apakah membutuhkan sumpit?"}
\end{convobox}

\begin{convobox}{38}
    \noindent
    \jpword{おお}{多}{Oo}\jpkana{くの}{ku no} 
    \jpkana{コンビニ}{konbini} \jpkana{では、}{de wa,} 
    \jpkana{お}{o}\jpword{はし}{箸}{hashi}\jpkana{、}{,} 
    \jpkana{フォーク、}{fooku,} 
    \jpkana{スプーンなどを}{supuun nado o} 
    \jpword{むりょう}{無料}{muryou} \jpkana{でつけてくれます。}{de tsukete kuremasu.}

    \vspace{1.5em}
    \noindent{\Large \color{articolor} \textbf{Arti:} Di banyak minimarket, mereka akan menyertakan sumpit, garpu, sendok, dll secara gratis.}
\end{convobox}

\begin{convobox}{39}
    \noindent
    \jpword{こた}{答}{Kota}\jpkana{えるときは、}{eru toki wa,} 
    \jpword{さき}{先}{saki}\jpkana{ほどと}{hodo to} 
    \jpword{おな}{同}{ona}\jpkana{じように}{ji you ni} 
    \jpword{い}{言}{i}\jpkana{えば}{eba} 
    \jpword{だいじょうぶ}{大丈夫}{daijoubu} \jpkana{です。}{desu.}

    \vspace{1.5em}
    \noindent{\Large \color{articolor} \textbf{Arti:} Saat menjawab, jika kamu mengatakannya sama seperti sebelumnya, itu sudah cukup.}
\end{convobox}

\begin{convobox}{40}
    \noindent
    \jpkana{「はい、}{"Hai,} \jpkana{お}{o}\jpword{ねが}{願}{nega}\jpkana{いします。」}{ishimasu."}

    \vspace{1.5em}
    \noindent{\Large \color{articolor} \textbf{Arti:} "Ya, tolong."}
\end{convobox}

\begin{convobox}{41}
    \noindent
    \jpkana{「いえ、}{"Ie,} \jpword{だいじょうぶ}{大丈夫}{daijoubu} \jpkana{です。」}{desu."}

    \vspace{1.5em}
    \noindent{\Large \color{articolor} \textbf{Arti:} "Tidak, tidak usah."}
\end{convobox}

\begin{convobox}{42}
    \noindent
    \jpkana{もし}{Moshi} \jpword{てんいん}{店員}{tenin}\jpkana{さんが}{san ga} 
    \jpkana{お}{o}\jpword{はし}{箸}{hashi}\jpkana{を}{o} 
    \jpkana{つけるのを}{tsukeru no o} 
    \jpword{わす}{忘}{wasu}\jpkana{れている}{rete iru} 
    \jpword{とき}{時}{toki}\jpkana{は、}{wa,} 
    \jpkana{このように}{kono you ni} 
    \jpword{い}{言}{i}\jpkana{いましょう。}{imashou.}

    \vspace{1.5em}
    \noindent{\Large \color{articolor} \textbf{Arti:} Jika kasir lupa memberikan sumpitnya, ucapkanlah seperti ini.}
\end{convobox}

\begin{convobox}{43}
    \noindent
    \jpkana{「お}{"O}\jpword{はし}{箸}{hashi} \jpkana{お}{o}\jpword{ねが}{願}{nega}\jpkana{いします。」}{ishimasu."}

    \vspace{1.5em}
    \noindent{\Large \color{articolor} \textbf{Arti:} "Tolong berikan sumpitnya."}
\end{convobox}

\begin{lessonbox}{Aksi Pahlawan Konbini: Detektif Alat Makan!}
    \textbf{Penjelasan:} Biasanya kasir akan menebak alat makan apa yang kamu butuhkan. Beli mie? Dapat \textbf{O-hashi} (Sumpit). Beli puding? Dapat \textbf{Supuun} (Sendok kecil).
    \\
    \textbf{🎬 ACTION / PRAKTEK:}
    Jika kamu membeli mie instan untuk dibawa pulang ke hotel, jangan lupa minta sumpitnya JIKA kasir tidak bertanya! 
    Angkat tangan sedikit, tatap kasir sambil tersenyum, dan ucapkan dengan jelas: \textbf{"O-hashi, onegaishimasu!"}.
\end{lessonbox}

% ================= TITLE: PHRASE 3 =================
\vspace{2em}
\begin{tcolorbox}[colback=boxframe, colframe=boxframe, rounded corners, boxrule=0pt, arc=3mm, left=2mm, top=2mm, bottom=2mm]
    \Large \bfseries \color{white} 🛍️ 4. Phrase 3 (Kantong Plastik)
\end{tcolorbox}
\vspace{1em}

\begin{convobox}{44}
    \noindent
    \jpkana{「}{"} \jpword{ふくろ}{袋}{Fukuro} \jpkana{は}{wa} 
    \jpword{りよう}{利用}{riyou} \jpkana{ですか？」}{desu ka?"}

    \vspace{1.5em}
    \noindent{\Large \color{articolor} \textbf{Arti:} "Apakah membutuhkan kantong plastik?"}
\end{convobox}

\begin{convobox}{45}
    \noindent
    \jpkana{「}{"}\jpword{ふくろ}{袋}{Fukuro}\jpkana{」は、}{" wa,} 
    \jpkana{「レジ}{"Reji }\jpword{ぶくろ}{袋}{bukuro}\jpkana{」と}{" to} 
    \jpword{い}{言}{i}\jpkana{うこともあります。}{u koto mo arimasu.}

    \vspace{1.5em}
    \noindent{\Large \color{articolor} \textbf{Arti:} Kata "Fukuro" kadang-kadang juga disebut "Reji-bukuro" (kantong belanja).}
\end{convobox}

\begin{convobox}{46}
    \noindent
    \jpword{ふくろ}{袋}{Fukuro} \jpkana{は}{wa} 
    \jpword{すうねんまえ}{数年前}{suunen-mae} \jpkana{まで}{made} 
    \jpword{むりょう}{無料}{muryou} \jpkana{でしたが、}{deshita ga,}

    \vspace{1.5em}
    \noindent{\Large \color{articolor} \textbf{Arti:} Sampai beberapa tahun lalu plastik itu gratis, tapi}
\end{convobox}

\begin{convobox}{47}
    \noindent
    \jpword{さいきん}{最近}{Saikin} \jpkana{は}{wa} 
    \jpword{ほとん}{殆}{hoton}\jpkana{どの}{do no} 
    \jpkana{コンビニ}{konbini} \jpkana{で}{de} 
    \jpword{ゆうりょう}{有料}{yuuryou} \jpkana{になりました。}{ni narimashita.}

    \vspace{1.5em}
    \noindent{\Large \color{articolor} \textbf{Arti:} Akhir-akhir ini, di hampir semua minimarket menjadi berbayar.}
\end{convobox}

\begin{convobox}{48}
    \noindent
    \jpword{ねだん}{値段}{Nedan} \jpkana{は}{wa} 
    \jpword{さんえん}{3円}{san-en} \jpkana{から}{kara} 
    \jpword{ごえん}{5円}{go-en} \jpkana{ぐらいです。}{gurai desu.}

    \vspace{1.5em}
    \noindent{\Large \color{articolor} \textbf{Arti:} Harganya kira-kira dari 3 yen sampai 5 yen.}
\end{convobox}

\begin{convobox}{49}
    \noindent
    \jpkana{この}{Kono} \jpword{しつもん}{質問}{shitsumon} \jpkana{にも、}{ni mo,} 
    \jpword{さき}{先}{saki}\jpkana{ほどと}{hodo to} 
    \jpword{おな}{同}{ona}\jpkana{じように}{ji you ni} 
    \jpword{こた}{答}{kota}\jpkana{えれば}{ereba} 
    \jpword{だいじょうぶ}{大丈夫}{daijoubu} \jpkana{です。}{desu.}

    \vspace{1.5em}
    \noindent{\Large \color{articolor} \textbf{Arti:} Pada pertanyaan ini juga, jika menjawab dengan cara yang sama seperti tadi sudah cukup.}
\end{convobox}

\begin{convobox}{50}
    \noindent
    \jpkana{「はい、}{"Hai,} \jpkana{お}{o}\jpword{ねが}{願}{nega}\jpkana{いします。」}{ishimasu."}

    \vspace{1.5em}
    \noindent{\Large \color{articolor} \textbf{Arti:} "Ya, tolong."}
\end{convobox}

\begin{convobox}{51}
    \noindent
    \jpkana{「いえ、}{"Ie,} \jpword{だいじょうぶ}{大丈夫}{daijoubu} \jpkana{です。」}{desu."}

    \vspace{1.5em}
    \noindent{\Large \color{articolor} \textbf{Arti:} "Tidak, tidak usah."}
\end{convobox}

\begin{lessonbox}{Aksi Pahlawan Konbini: Siapkan Tas Belanjamu!}
    \textbf{Penjelasan:} Kata kunci yang harus kamu dengar dari kasir adalah \textbf{FUKURO} (Kantong plastik) atau \textbf{REJI-BUKURO}.
    \\
    \textbf{🎬 ACTION / PRAKTEK:} 
    Jadilah pembeli yang ramah lingkungan! Saat kasir menyebut kata \textit{Fukuro}, angkat tas kain yang kamu bawa dari rumah, senyum, dan katakan: \textbf{"Ie, daijoubu desu!"}. Kamu hemat 3 yen dan menyelamatkan bumi! 
\end{lessonbox}

% ================= TITLE: PHRASE 4 =================
\vspace{2em}
\begin{tcolorbox}[colback=boxframe, colframe=boxframe, rounded corners, boxrule=0pt, arc=3mm, left=2mm, top=2mm, bottom=2mm]
    \Large \bfseries \color{white} 💳 5. Phrase 4 (Kartu Poin)
\end{tcolorbox}
\vspace{1em}

\begin{convobox}{52}
    \noindent
    \jpkana{「ポイントカード}{"Pointo kaado} \jpkana{は}{wa} 
    \jpkana{お}{o}\jpword{も}{持}{mo}\jpkana{ちですか？」}{chi desu ka?"}

    \vspace{1.5em}
    \noindent{\Large \color{articolor} \textbf{Arti:} "Apakah Anda membawa kartu poin?"}
\end{convobox}

\begin{convobox}{53}
    \noindent
    \jpkana{コンビニ}{Konbini} \jpkana{には、}{ni wa,} 
    \jpkana{それぞれ}{sorezore} 
    \jpkana{ポイントカード}{pointo kaado} \jpkana{があります。}{ga arimasu.}

    \vspace{1.5em}
    \noindent{\Large \color{articolor} \textbf{Arti:} Di minimarket, masing-masing memiliki kartu poinnya sendiri.}
\end{convobox}

\begin{convobox}{54}
    \noindent
    \jpkana{お}{O}\jpword{かいけい}{会計}{kaikei} \jpkana{の}{no} 
    \jpword{まえ}{前}{mae} \jpkana{に、}{ni,} 
    \jpkana{このように}{kono you ni} 
    \jpword{き}{聞}{ki}\jpkana{かれることが}{kareru koto ga} 
    \jpword{おお}{多}{oo}\jpkana{いです。}{i desu.}

    \vspace{1.5em}
    \noindent{\Large \color{articolor} \textbf{Arti:} Sebelum melakukan pembayaran, kamu sering ditanya seperti ini.}
\end{convobox}

\begin{convobox}{55}
    \noindent
    \jpkana{ポイントカード}{Pointo kaado} \jpkana{を}{o} 
    \jpword{も}{持}{mo}\jpkana{っている}{tte iru} 
    \jpword{とき}{時}{toki}\jpkana{は、}{wa,} 
    \jpkana{「はい」と}{"Hai" to} \jpword{い}{言}{i}\jpkana{って}{tte}

    \vspace{1.5em}
    \noindent{\Large \color{articolor} \textbf{Arti:} Saat kamu membawa/memiliki kartu poin, katakan "Ya", lalu}
\end{convobox}

\begin{convobox}{56}
    \noindent
    \jpkana{カード}{Kaado} \jpkana{や}{ya} \jpkana{アプリの}{apuri no} 
    \jpkana{バーコードを}{baakoodo o} 
    \jpword{てんいん}{店員}{tenin}\jpkana{さんに}{san ni} 
    \jpword{み}{見}{mi}\jpkana{せます。}{semasu.}

    \vspace{1.5em}
    \noindent{\Large \color{articolor} \textbf{Arti:} tunjukkan kartu atau barcode aplikasi kepada kasir.}
\end{convobox}

\begin{convobox}{57}
    \noindent
    \jpword{も}{持}{Mo}\jpkana{っていない}{tte inai} 
    \jpword{とき}{時}{toki}\jpkana{は、}{wa,} 
    \jpkana{「いえ。」と}{"Ie." to} 
    \jpword{い}{言}{i}\jpkana{えば}{eba} 
    \jpword{だいじょうぶ}{大丈夫}{daijoubu} \jpkana{です。}{desu.}

    \vspace{1.5em}
    \noindent{\Large \color{articolor} \textbf{Arti:} Saat tidak punya, bilang "Tidak" saja sudah cukup.}
\end{convobox}

\begin{lessonbox}{Aksi Pahlawan Konbini: Turis Tidak Butuh Poin!}
    \textbf{Penjelasan:} Kartu poin biasanya dikumpulkan oleh orang lokal Jepang yang sering belanja ke sana (seperti T-Point atau d-Point).
    \\
    \textbf{🎬 ACTION / PRAKTEK:} 
    Karena kamu sedang jalan-jalan (turis), kamu tidak butuh ini! Jika kasir bilang kata \textbf{POINTO KAADO}, segera gelengkan kepalamu dan bilang dengan tegas tapi sopan: \textbf{"Ie!"} atau \textbf{"Nai desu"} (Tidak ada).
\end{lessonbox}

% ================= TITLE: PHRASE 5 =================
\vspace{2em}
\begin{tcolorbox}[colback=boxframe, colframe=boxframe, rounded corners, boxrule=0pt, arc=3mm, left=2mm, top=2mm, bottom=2mm]
    \Large \bfseries \color{white} 📱 6. Phrase 5 (Metode Pembayaran)
\end{tcolorbox}
\vspace{1em}

\begin{convobox}{58}
    \noindent
    \jpkana{「Tanaka Pay}{"Tanaka Pay} \jpkana{で}{de} 
    \jpkana{お}{o}\jpword{ねが}{願}{nega}\jpkana{いします。」}{ishimasu."}

    \vspace{1.5em}
    \noindent{\Large \color{articolor} \textbf{Arti:} "Pakai Tanaka Pay, tolong."}
\end{convobox}

\begin{convobox}{59}
    \noindent
    \jpkana{お}{O}\jpword{かいけい}{会計}{kaikei} \jpkana{の}{no} 
    \jpword{ほうほう}{方法}{houhou} \jpkana{を}{o} 
    \jpword{つた}{伝}{tsuta}\jpkana{える}{eru} 
    \jpkana{フレーズです。}{fureezu desu.}

    \vspace{1.5em}
    \noindent{\Large \color{articolor} \textbf{Arti:} Ini adalah frasa untuk menyampaikan metode pembayaran.}
\end{convobox}

\begin{convobox}{60}
    \noindent
    \jpkana{コンビニ}{Konbini} \jpkana{では、}{de wa,} 
    \jpword{げんきん}{現金}{genkin}\jpkana{、}{,} 
    \jpkana{クレジットカード、}{kurejitto kaado,} 
    \jpkana{スマホ}{sumaho}\jpword{けっさい}{決済}{kessai} \jpkana{など、}{nado,}

    \vspace{1.5em}
    \noindent{\Large \color{articolor} \textbf{Arti:} Di minimarket, ada uang tunai (Genkin), kartu kredit, pembayaran smartphone, dll,}
\end{convobox}

\begin{convobox}{61}
    \noindent
    \jpkana{いろいろな}{Iroiro na} 
    \jpword{ほうほう}{方法}{houhou} \jpkana{で}{de} 
    \jpword{しはら}{支払}{shihara}\jpkana{いができます。}{i ga dekimasu.}

    \vspace{1.5em}
    \noindent{\Large \color{articolor} \textbf{Arti:} kamu bisa membayar dengan berbagai macam metode.}
\end{convobox}

\begin{convobox}{62}
    \noindent
    \jpword{たと}{例}{Tato}\jpkana{えば}{eba} 
    \jpkana{クレジットカードで}{kurejitto kaado de} 
    \jpword{はら}{払}{hara}\jpkana{いたい}{itai} 
    \jpword{とき}{時}{toki}\jpkana{は、}{wa,} 
    \jpkana{このように}{kono you ni} 
    \jpword{い}{言}{i}\jpkana{います。}{imasu.}

    \vspace{1.5em}
    \noindent{\Large \color{articolor} \textbf{Arti:} Misalnya, saat ingin membayar dengan kartu kredit, ucapkanlah seperti ini.}
\end{convobox}

\begin{convobox}{63}
    \noindent
    \jpkana{「クレジットカードで}{"Kurejitto kaado de} 
    \jpkana{お}{o}\jpword{ねが}{願}{nega}\jpkana{いします。」}{ishimasu."}

    \vspace{1.5em}
    \noindent{\Large \color{articolor} \textbf{Arti:} "Pakai Kartu Kredit, tolong."}
\end{convobox}

\begin{lessonbox}{Aksi Pahlawan Konbini: Sebut Senjatamu!}
    \textbf{Penjelasan:} Di Jepang, mesin kasir harus ditekan sesuai alat bayarmu. Jadi kamu HARUS kasih tahu kasir mau bayar pakai apa!
    \\
    \textbf{🎬 ACTION / PRAKTEK:} 
    Pegang alat bayarmu (Uang cash, Kartu, atau HP). Tunjukkan pada kasir, dan gunakan rumus sakti ini: \textbf{[Alat Bayar] + de onegaishimasu!}
    \begin{itemize}
        \item Uang Tunai: \textbf{"Genkin de onegaishimasu!"}
        \item Kartu IC (Suica/Pasmo): \textbf{"Suica de onegaishimasu!"}
        \item Kartu Kredit: \textbf{"Kaado de onegaishimasu!"}
    \end{itemize}
\end{lessonbox}

% ================= TITLE: PHRASE 6 =================
\vspace{2em}
\begin{tcolorbox}[colback=boxframe, colframe=boxframe, rounded corners, boxrule=0pt, arc=3mm, left=2mm, top=2mm, bottom=2mm]
    \Large \bfseries \color{white} 🧾 7. Phrase 6 (Struk Belanja)
\end{tcolorbox}
\vspace{1em}

\begin{convobox}{64}
    \noindent
    \jpkana{「レシート}{"Reshiito} \jpword{りよう}{利用}{riyou} \jpkana{ですか？」}{desu ka?"}

    \vspace{1.5em}
    \noindent{\Large \color{articolor} \textbf{Arti:} "Apakah butuh struk?"}
\end{convobox}

\begin{convobox}{65}
    \noindent
    \jpkana{お}{O}\jpword{かいけい}{会計}{kaikei} \jpkana{の}{no} 
    \jpword{あと}{後}{ato} \jpkana{に、}{ni,} 
    \jpkana{よく}{yoku} 
    \jpword{き}{聞}{ki}\jpkana{かれる}{kareru} 
    \jpkana{フレーズです。}{fureezu desu.}

    \vspace{1.5em}
    \noindent{\Large \color{articolor} \textbf{Arti:} Ini adalah frasa yang sering ditanyakan setelah pembayaran selesai.}
\end{convobox}

\begin{convobox}{66}
    \noindent
    \jpkana{これも}{Kore mo} 
    \jpkana{さっきと}{sakki to} 
    \jpword{おな}{同}{ona}\jpkana{じように}{ji you ni} 
    \jpword{こた}{答}{kota}\jpkana{えれば}{ereba} 
    \jpword{だいじょうぶ}{大丈夫}{daijoubu} \jpkana{です。}{desu.}

    \vspace{1.5em}
    \noindent{\Large \color{articolor} \textbf{Arti:} Ini juga jika dijawab sama seperti tadi sudah cukup kok.}
\end{convobox}

\begin{convobox}{67}
    \noindent
    \jpkana{「はい、}{"Hai,} \jpkana{お}{o}\jpword{ねが}{願}{nega}\jpkana{いします。」}{ishimasu."}

    \vspace{1.5em}
    \noindent{\Large \color{articolor} \textbf{Arti:} "Ya, tolong."}
\end{convobox}

\begin{convobox}{68}
    \noindent
    \jpkana{「いえ、}{"Ie,} \jpword{だいじょうぶ}{大丈夫}{daijoubu} \jpkana{です。」}{desu."}

    \vspace{1.5em}
    \noindent{\Large \color{articolor} \textbf{Arti:} "Tidak, tidak usah."}
\end{convobox}

% ================= TITLE: TIPS =================
\vspace{2em}
\begin{tcolorbox}[colback=boxframe, colframe=boxframe, rounded corners, boxrule=0pt, arc=3mm, left=2mm, top=2mm, bottom=2mm]
    \Large \bfseries \color{white} 💡 8. Tips Rahasia Konbini (Toilet)
\end{tcolorbox}
\vspace{1em}

\begin{convobox}{69}
    \noindent
    \jpword{おお}{多}{Oo}\jpkana{くの}{ku no} 
    \jpkana{コンビニ}{konbini} \jpkana{には、}{ni wa,} 
    \jpkana{トイレが}{toire ga} 
    \jpkana{あります。}{arimasu.}

    \vspace{1.5em}
    \noindent{\Large \color{articolor} \textbf{Arti:} Di banyak minimarket, terdapat toilet.}
\end{convobox}

\begin{convobox}{70}
    \noindent
    \jpkana{トイレを}{Toire o} 
    \jpword{か}{借}{ka}\jpkana{りるときは、}{riru toki wa,} 
    \jpword{てんいん}{店員}{tenin}\jpkana{さんに}{san ni} 
    \jpword{ひとこえ}{一声}{hito-koe} 
    \jpkana{かけましょう。}{kakemashou.}

    \vspace{1.5em}
    \noindent{\Large \color{articolor} \textbf{Arti:} Saat meminjam (menggunakan) toilet, mari menyapa/meminta izin kasir terlebih dahulu.}
\end{convobox}

\begin{convobox}{71}
    \noindent
    \jpkana{「すみません、}{"Sumimasen,} 
    \jpkana{トイレを}{toire o} 
    \jpword{か}{借}{ka}\jpkana{りてもいいですか？」}{rite mo ii desu ka?"}

    \vspace{1.5em}
    \noindent{\Large \color{articolor} \textbf{Arti:} "Maaf (Permisi), bolehkah saya meminjam toilet?"}
\end{convobox}

\begin{convobox}{72}
    \noindent
    \jpkana{トイレの}{Toire no} 
    \jpword{りよう}{利用}{riyou} \jpkana{は}{wa} 
    \jpword{むりょう}{無料}{muryou} \jpkana{ですが、}{desu ga,}

    \vspace{1.5em}
    \noindent{\Large \color{articolor} \textbf{Arti:} Penggunaan toilet itu gratis, tapi}
\end{convobox}

\begin{convobox}{73}
    \noindent
    \jpkana{できたら}{dekitara} 
    \jpword{なに}{何}{nani}\jpkana{か}{ka} 
    \jpword{しょうひん}{商品}{shouhin} \jpkana{を}{o} 
    \jpword{か}{買}{ka}\jpkana{って}{tte} 
    \jpword{みせ}{店}{mise} \jpkana{を}{o} 
    \jpword{で}{出}{de}\jpkana{ると}{ru to} 
    \jpkana{いいですね。}{ii desu ne.}

    \vspace{1.5em}
    \noindent{\Large \color{articolor} \textbf{Arti:} jika memungkinkan, lebih baik membeli suatu produk (apa saja) sebelum keluar dari toko ya.}
\end{convobox}

\begin{lessonbox}{Aksi Pahlawan Konbini: Sopan Santun Toilet!}
    \textbf{Penjelasan:} Di Jepang, toilet Konbini adalah penyelamat turis! Tapi ingat tata kramanya ya.
    \\
    \textbf{🎬 ACTION / PRAKTEK:} 
    Jangan langsung nyelonong masuk! Hampiri kasir, katakan \textbf{"Toire o karite mo ii desu ka?"}. Setelah selesai buang air, jangan langsung pergi. Ambil air minum botol seharga 100 yen atau permen kecil, lalu bayar di kasir sebagai rasa terima kasih. Ini adalah etika yang sangat keren!
\end{lessonbox}

% ================= TITLE: ROLE PLAY =================
\vspace{2em}
\begin{tcolorbox}[colback=boxframe, colframe=boxframe, rounded corners, boxrule=0pt, arc=3mm, left=2mm, top=2mm, bottom=2mm]
    \Large \bfseries \color{white} 🗣️ 9. Waktunya Role Play! (Praktik Ucapkan)
\end{tcolorbox}
\vspace{1em}

\begin{convobox}{74}
    \noindent
    \jpkana{それでは、}{Sore dewa,} 
    \jpword{いま}{今}{ima} \jpkana{から}{kara} 
    \jpkana{ロールプレイを}{rooru purei o} \jpkana{します。}{shimasu.}

    \vspace{1.5em}
    \noindent{\Large \color{articolor} \textbf{Arti:} Kalau begitu, sekarang kita akan melakukan roleplay (bermain peran).}
\end{convobox}

\begin{convobox}{75}
    \noindent
    \jpkana{こんかいは、}{Konkai wa,} 
    \jpkana{みなさんは}{minasan wa} 
    \jpkana{お}{o}\jpword{きゃく}{客}{kyaku}\jpkana{さんに}{san ni} 
    \jpkana{なったつもりで、}{natta tsumori de,} 
    \jpword{こえ}{声}{koe} \jpkana{に}{ni} 
    \jpword{だ}{出}{da}\jpkana{して}{shite} 
    \jpword{はな}{話}{hana}\jpkana{してみてください。}{shite mite kudasai.}

    \vspace{1.5em}
    \noindent{\Large \color{articolor} \textbf{Arti:} Kali ini, bayangkanlah dirimu sebagai pembeli, dan cobalah berbicara dengan suara yang lantang.}
\end{convobox}

\begin{lessonbox}{⚠️ TUGAS TERAKHIR PAHLAWAN KONBINI! ⚠️}
    Anggap saja kamu sedang berdiri di depan kasir Jepang membawa kotak makan siang. Baca keras-keras Romaji di bagian bawah (yang menjadi jawabanmu). \textbf{Jangan dibaca dalam hati!} Mulutmu harus terbiasa mengucapkan kata-kata ini agar lancar saat di Jepang nanti. Ayo kita mulai!
\end{lessonbox}

\begin{convobox}{76}
    \noindent
    \jpkana{【店員】いらっしゃいませ。}{[Tenin] Irasshaimase.} 
    \jpkana{こちら}{Kochira} \jpword{あたた}{温}{atata}\jpkana{めますか？}{memasu ka?}

    \vspace{1.5em}
    \noindent{\Large \color{articolor} \textbf{Arti:} [Kasir] Selamat datang. Apakah mau dipanaskan?}
\end{convobox}

\begin{convobox}{77}
    \noindent
    \jpkana{【あなた】}{[Anata]} \jpkana{はい、}{Hai,} 
    \jpkana{お}{o}\jpword{ねが}{願}{nega}\jpkana{いします。}{ishimasu.}

    \vspace{1.5em}
    \noindent{\Large \color{articolor} \textbf{Arti:} [Kamu] Ya, tolong.}
\end{convobox}

\begin{convobox}{78}
    \noindent
    \jpkana{【店員】お}{[Tenin] O}\jpword{はし}{箸}{hashi} \jpkana{は}{wa} 
    \jpword{りよう}{利用}{riyou} \jpkana{ですか？}{desu ka?}

    \vspace{1.5em}
    \noindent{\Large \color{articolor} \textbf{Arti:} [Kasir] Pakai sumpit?}
\end{convobox}

\begin{convobox}{79}
    \noindent
    \jpkana{【あなた】}{[Anata]} \jpkana{はい、}{Hai,} 
    \jpkana{お}{o}\jpword{ねが}{願}{nega}\jpkana{いします。}{ishimasu.}

    \vspace{1.5em}
    \noindent{\Large \color{articolor} \textbf{Arti:} [Kamu] Ya, tolong.}
\end{convobox}

\begin{convobox}{80}
    \noindent
    \jpkana{【店員】}{[Tenin] }\jpword{ふくろ}{袋}{Fukuro} \jpkana{は}{wa} 
    \jpword{りよう}{利用}{riyou} \jpkana{ですか？}{desu ka?}

    \vspace{1.5em}
    \noindent{\Large \color{articolor} \textbf{Arti:} [Kasir] Butuh kantong plastik?}
\end{convobox}

\begin{convobox}{81}
    \noindent
    \jpkana{【あなた】}{[Anata]} \jpkana{いえ、}{Ie,} \jpword{だいじょうぶ}{大丈夫}{daijoubu} \jpkana{です。}{desu.}

    \vspace{1.5em}
    \noindent{\Large \color{articolor} \textbf{Arti:} [Kamu] (Aku bawa tas kain, jadi) Tidak usah.}
\end{convobox}

\begin{convobox}{82}
    \noindent
    \jpkana{【店員】お}{[Tenin] O}\jpword{かいけい}{会計}{kaikei} 
    \jpword{ろっぴゃくえん}{600円}{roppyaku-en} 
    \jpkana{でございます。}{de gozaimasu.}

    \vspace{1.5em}
    \noindent{\Large \color{articolor} \textbf{Arti:} [Kasir] Totalnya 600 yen.}
\end{convobox}

\begin{convobox}{83}
    \noindent
    \jpkana{【あなた】クレジットカードで}{[Anata] Kurejitto kaado de} 
    \jpkana{お}{o}\jpword{ねが}{願}{nega}\jpkana{いします。}{ishimasu.}

    \vspace{1.5em}
    \noindent{\Large \color{articolor} \textbf{Arti:} [Kamu] Pakai Kartu Kredit, tolong.}
\end{convobox}

\begin{convobox}{84}
    \noindent
    \jpkana{【店員】かしこまりました。}{[Tenin] Kashikomarimashita.} 
    \jpkana{レシート}{Reshiito} \jpword{りよう}{利用}{riyou} \jpkana{ですか？}{desu ka?}

    \vspace{1.5em}
    \noindent{\Large \color{articolor} \textbf{Arti:} [Kasir] Baik. Butuh struk?}
\end{convobox}

\begin{convobox}{85}
    \noindent
    \jpkana{【あなた】}{[Anata]} \jpkana{はい、}{Hai,} 
    \jpkana{お}{o}\jpword{ねが}{願}{nega}\jpkana{いします。}{ishimasu.}

    \vspace{1.5em}
    \noindent{\Large \color{articolor} \textbf{Arti:} [Kamu] Ya, tolong (buat kenang-kenangan!).}
\end{convobox}

\begin{convobox}{86}
    \noindent
    \jpkana{【店員】ありがとうございました。}{[Tenin] Arigatou gozaimashita.} 
    \jpkana{また}{Mata} \jpkana{お}{o}\jpword{こ}{越}{ko}\jpkana{しくださいませ。}{shi kudasaimase.}

    \vspace{1.5em}
    \noindent{\Large \color{articolor} \textbf{Arti:} [Kasir] Terima kasih banyak. Silakan datang kembali.}
\end{convobox}

% --- SUMMARY BOX ---
\vspace{2em}
\begin{tcolorbox}[
    colback=blue!5!white,
    colframe=blue!80!black,
    boxrule=3pt,
    arc=5mm,
    title={\huge\bfseries \sffamily 🌟 Rangkuman Misi Konbini 🌟},
    fonttitle=\color{white},
    attach boxed title to top center={yshift=-4mm},
    boxed title style={colback=blue!80!black, rounded corners},
    left=5mm, right=5mm, top=7mm, bottom=5mm
]
\Large \textbf{Sempurna! Kamu telah lulus ujian Konbini Jepang!}
\begin{itemize}
    \item Ingat 2 Jawaban Sakti: \textbf{Hai, onegaishimasu} (Ya tolong) dan \textbf{Ie, daijoubu desu} (Tidak usah).
    \item Dengarkan kata kunci: \textit{Atatamemasu ka} (Panaskan?), \textit{O-hashi} (Sumpit), \textit{Fukuro} (Plastik), \textit{Pointo Kaado} (Kartu poin), dan \textit{Reshiito} (Struk).
    \item Cara bayar: Sebutkan nama alat bayarnya + \textit{De onegaishimasu}.
    \item Toilet darurat: Selalu sapa kasir \textit{Toire o karitemo ii desu ka?} dan belilah barang kecil sebagai bentuk kesopanan.
\end{itemize}
Sekarang kamu sudah 100\% siap berbelanja di minimarket Jepang tanpa perlu merasa panik. Terus semangat belajarnya, dan \textit{Mata jikai no douga de o-ai shimashou!} (Sampai jumpa di video/materi selanjutnya!)
\end{tcolorbox}

\newpage